\section{Technical Overview}

The proof system is best viewed as a certification procedure for quantized GBDT
training. The main challenge is not tree inference itself, but proving that
every boosting round computes the correct split statistics and commits a split
that optimizes the training objective. Two obstacles recur throughout the
construction. First, the proof naturally produces claims over different
polynomial domains and tensor layouts. Second, plaintext split search is
usually implemented through sorting or feature scans, which are awkward proof
targets. The overview below therefore isolates three ingredients: the
committed forest interface, the algebraic rewriting of split search, and the
batching tools that make the resulting checks compose succinctly.

\subsection{Forest Representation and Batched Inference}

We adopt the hyperbox forest encoding of \cite{cq++24}. Each node carries a
feature-wise half-open interval box, and a sample reaches a leaf exactly when
its quantized feature vector lies inside that box. All committed objects share
one global node namespace, so the row index
$u \in [N_{\mathrm{tot}}]$ refers to the same node everywhere. The committed
forest consists of:
\begin{itemize}
  \item routing matrices
  $\mathbf{L},\mathbf{R}\in\{0,1\}^{N_{\mathrm{tot}}\times N_{\mathrm{tot}}}$,
  where an internal-node row is one-hot and selects the left or
  right child of that node, while leaf rows are zero;
  \item lower and upper node bounds
  $\mathbf{N}^{\leftarrow},\mathbf{N}^{\rightarrow}\in[B+1]^{N_{\mathrm{tot}}\times d}$,
  where row $u$ stores the $d$-dimensional half-open box of node $u$ and every
  leaf row describes one accepting region of the tree;
  \item a feature-selection matrix
  $\mathbf{E}\in\{0,1\}^{N_{\mathrm{tot}}\times d}$,
  whose internal-node rows are one-hot and indicate which feature is split at
  each node, while leaf rows are zero;
  \item a threshold-index vector
  $\boldsymbol{\Theta}\in[B]^{N_{\mathrm{tot}}}$,
  whose internal-node entries record the selected threshold bin and whose leaf
  and padding rows are zero;
  \item a tree-index vector $\mathbf{tree}\in[TK]^{N_{\mathrm{tot}}}$ and a
  root matrix $\mathbf{Root}\in\{0,1\}^{TK\times N_{\mathrm{tot}}}$;
  \item a leaf-score vector $\mathbf{score}\in\mathbb{F}^{N_{\mathrm{tot}}}$,
  whose leaf entries store the scalar contribution returned by that tree and
  whose internal-node entries are zero or padding.
\end{itemize}
The pair $(\mathbf{E},\boldsymbol{\Theta})$ witnesses that
$\mathbf{N}^{\leftarrow}$ and $\mathbf{N}^{\rightarrow}$ encode a valid
hyperbox forest. The routing matrices $(\mathbf{L},\mathbf{R})$ are reused in
training to propagate histogram statistics from leaves to internal nodes.

\Cref{fig:forest-encoding-inference} follows the exposition style of
\cite{cq++24}: the left panel presents two concrete trees in the global node
namespace, while the right panel shows the committed row view on which the
proof operates. The committed columns record the tree index, split metadata,
node box, and leaf score. The online fetch opens from this view only the tree
index, node box, and score of the accepting leaf.

\begin{figure*}[t]
  \centering
  \resizebox{\linewidth}{!}{%
  \begin{tikzpicture}[
    >=Latex,
    font=\small,
    paneltitle/.style={font=\bfseries},
    treelabel/.style={font=\bfseries\small},
    internal/.style={
      circle,
      draw=blue!55!black,
      line width=0.55pt,
      fill=blue!5,
      minimum size=10.2mm,
      inner sep=1pt,
      align=center
    },
    leaf/.style={
      rounded corners=2pt,
      draw=green!50!black,
      line width=0.55pt,
      fill=green!10,
      minimum width=14mm,
      minimum height=8.6mm,
      inner sep=2pt,
      align=center
    },
    note/.style={align=center, font=\scriptsize},
    bigtable/.style={
      matrix of nodes,
      ampersand replacement=\&,
      row sep=0.75mm,
      column sep=\colgap,
      nodes={
        anchor=center,
        inner xsep=2.5pt,
        inner ysep=1.2pt,
        minimum height=4.5mm,
        text depth=.25ex,
        outer sep=0pt
      },
      row 1/.style={nodes={font=\bfseries}},
      column 1/.style={nodes={minimum width=7mm}},
      column 2/.style={nodes={minimum width=11mm}},
      column 3/.style={nodes={minimum width=10mm}},
      column 4/.style={nodes={minimum width=13mm}},
      column 5/.style={nodes={minimum width=9mm}},
      column 6/.style={nodes={minimum width=13mm}},
      column 7/.style={nodes={minimum width=13mm}},
      column 8/.style={nodes={minimum width=12mm}}
    },
    colbg/.style={
      draw=gray!70,
      rounded corners=2pt,
      line width=.5pt,
      fill=white,
      inner sep=2.8pt
    },
    meta/.style={
      draw=gray!70,
      rounded corners=2pt,
      fill=gray!6,
      align=center,
      inner sep=6pt
    },
    fetchrow/.style={
      draw=red!75!black,
      rounded corners=2pt,
      line width=.85pt,
      inner xsep=4pt,
      inner ysep=2pt
    },
    fetchtxt/.style={text=red!75!black, font=\bfseries},
    outer/.style={
      draw=gray!55,
      dashed,
      rounded corners=4pt,
      inner sep=8pt
    }
  ]

    \def\colgap{4.2mm}

    % =========================
    % Left panel: tree semantics
    % =========================
    \node[paneltitle] (ltitle) at (-4.8,3.12) {(a) Tree semantics};

    % Tree 1
    \node[treelabel] at (-7.55,2.42) {tree $1$};
    \node[internal] (t0r) at (-7.55,1.72) {$1$\\[-1mm]\scriptsize$x_1<3$};
    \node[internal] (t0l) at (-9.05,0.38) {$2$\\[-1mm]\scriptsize$x_2<4$};
    \node[leaf]     (t0rleaf) at (-6.00,0.38) {$7$\\[-1mm]\scriptsize$w_2$};
    \node[leaf]     (t0w0) at (-10.00,-1.02) {$5$\\[-1mm]\scriptsize$w_0$};
    \node[leaf]     (t0w1) at (-8.35,-1.02) {$6$\\[-1mm]\scriptsize$w_1$};

    \draw[->, line width=.75pt] (t0r) -- (t0l);
    \draw[->, line width=.75pt] (t0r) -- (t0rleaf);
    \draw[->, line width=.75pt] (t0l) -- (t0w0);
    \draw[->, line width=.75pt] (t0l) -- (t0w1);

    % Tree 2
    \node[treelabel] at (-2.65,2.42) {tree $2$};
    \node[internal] (t1r) at (-2.65,1.72) {$3$\\[-1mm]\scriptsize$x_1<2$};
    \node[leaf]     (t1lleaf) at (-4.15,0.38) {$8$\\[-1mm]\scriptsize$w'_0$};
    \node[internal] (t1rr) at (-1.15,0.38) {$4$\\[-1mm]\scriptsize$x_2<6$};
    \node[leaf]     (t1w1) at (-2.05,-1.02) {$9$\\[-1mm]\scriptsize$w'_1$};
    \node[leaf]     (t1w2) at (-0.40,-1.02) {$10$\\[-1mm]\scriptsize$w'_2$};

    \draw[->, line width=.75pt] (t1r) -- (t1lleaf);
    \draw[->, line width=.75pt] (t1r) -- (t1rr);
    \draw[->, line width=.75pt] (t1rr) -- (t1w1);
    \draw[->, line width=.75pt] (t1rr) -- (t1w2);

    % =====================================
    % Right panel: one unified committed table
    % =====================================
    \node[paneltitle] (rtitle) at (7.25,3.12) {(b) Committed row view used by the proof};

    \matrix[bigtable, anchor=north west] (tbl) at (1.15,2.55) {
      $\mathbf{u}$ \& $\mathbf{tree}$ \& $\mathbf{is\_leaf}$ \& $\mathbf{E}$ \& $\boldsymbol{\Theta}$ \& $\mathbf{N}^{\leftarrow}$ \& $\mathbf{N}^{\rightarrow}$ \& $\mathbf{score}$ \\
      $0$ \& $0$ \& $0$ \& $(0,0)$ \& $0$ \& $(0,0)$ \& $(0,0)$ \& $0$ \\
      $1$ \& $1$ \& $0$ \& $(1,0)$ \& $3$ \& $(0,0)$ \& $(B,B)$ \& $0$ \\
      $2$ \& $1$ \& $0$ \& $(0,1)$ \& $4$ \& $(0,0)$ \& $(3,B)$ \& $0$ \\
      $3$ \& $2$ \& $0$ \& $(1,0)$ \& $2$ \& $(0,0)$ \& $(B,B)$ \& $0$ \\
      $4$ \& $2$ \& $0$ \& $(0,1)$ \& $6$ \& $(2,0)$ \& $(B,B)$ \& $0$ \\
      $5$ \& $1$ \& $1$ \& $(0,0)$ \& $0$ \& $(0,0)$ \& $(3,4)$ \& $w_0$ \\
      $6$ \& $1$ \& $1$ \& $(0,0)$ \& $0$ \& $(0,4)$ \& $(3,B)$ \& $w_1$ \\
      |[fetchtxt]| $7$ \& |[fetchtxt]| $1$ \& |[fetchtxt]| $1$ \& |[fetchtxt]| $(0,0)$ \& |[fetchtxt]| $0$ \& |[fetchtxt]| $(3,0)$ \& |[fetchtxt]| $(B,B)$ \& |[fetchtxt]| $w_2$ \\
      $8$ \& $2$ \& $1$ \& $(0,0)$ \& $0$ \& $(0,0)$ \& $(2,B)$ \& $w'_0$ \\
      |[fetchtxt]| $9$ \& |[fetchtxt]| $2$ \& |[fetchtxt]| $1$ \& |[fetchtxt]| $(0,0)$ \& |[fetchtxt]| $0$
        \& |[fetchtxt]| $(2,0)$ \& |[fetchtxt]| $(B,6)$ \& |[fetchtxt]| $w'_1$ \\
      $10$ \& $2$ \& $1$ \& $(0,0)$ \& $0$ \& $(2,6)$ \& $(B,B)$ \& $w'_2$ \\
    };

    % Column background boxes
    \begin{scope}[on background layer]
      \foreach \c in {1,...,8}{
        \node[colbg, fit=(tbl-1-\c)(tbl-12-\c)] {};
      }
    \end{scope}

    % Horizontal separators: after header and after internal rows
    \foreach \r in {1,6}{
      \draw[gray!60, line width=.4pt]
        ([xshift=-2pt]tbl-\r-1.south west) -- ([xshift=2pt]tbl-\r-8.south east);
    }

    % Highlight fetched rows (one row per tree)
    \node[fetchrow, fit=(tbl-9-1)(tbl-9-8)] (fetchA) {};
    \node[fetchrow, fit=(tbl-11-1)(tbl-11-8)] (fetchB) {};

    % Query box
    \node[meta, anchor=north west, text width=3.35cm] (sample)
      at ([xshift=6mm]tbl.north east)
      {query $\mathbf{x}=(4,2)$\\[1mm]
       fetch one row per tree:\\
        $u_1=7,\;u_2=9$\\[1mm]
       and prove, for any fetched row $u$,\\
       $\mathbf{N}^{\leftarrow}[u]\le \mathbf{x}<\mathbf{N}^{\rightarrow}[u]$};

    % Arrows to the two fetched rows
    \draw[->, very thick, red!75!black]
      ([yshift=7pt]sample.west) to[out=180,in=0] ([xshift=2pt]fetchA.east);
    \draw[->, very thick, red!75!black]
      ([yshift=-7pt]sample.west) to[out=180,in=0] ([xshift=2pt]fetchB.east);

    % Outer frame
    \node[
      outer,
      fit=(ltitle)(rtitle)
          (t0r)(t0l)(t0rleaf)(t0w0)(t0w1)
          (t1r)(t1lleaf)(t1rr)(t1w1)(t1w2)
          (tbl)(sample)
    ] {};

  \end{tikzpicture}%
  }
  \caption{An illustration of the forest encoding and inference semantics.
  The left panel shows two compact trees with
  global node identifiers, and the right panel shows the committed
  objects used by inference. The highlighted entries are the fetched rows for one
  query, one per tree. Batched inference performs this fetch for every
  tree-sample pair in one shot and then aggregates scores across trees for
  each sample-class pair.}
  \Description{Two example decision trees and their committed table
  representation. The table includes node identifiers, tree identifiers,
  leaf indicators, split-feature masks, thresholds, hyper-box bounds, and leaf
  scores. Two leaf rows are highlighted to show how inference fetches one row
  per tree and checks interval containment for a query point.}
  \label{fig:forest-encoding-inference}
\end{figure*}


This encoding supports both inference and training. For inference, the protocol
follows the row-fetch pattern of \cite{cq++24}. Given a query batch
$\mathbf{X}\in[B]^{b\times d}$, the prover conceptually replicates
$\mathbf{X}$ once per tree, determines the accepting leaf row of each
tree-sample pair, and proves a batched fetch of
\[
(\mathbf{tree},\mathbf{N}^{\leftarrow},\mathbf{N}^{\rightarrow},\mathbf{score})
\]
from those committed leaf rows. It then proves that the replicated query lies
inside the fetched half-open box
$[\mathbf{N}^{\leftarrow},\mathbf{N}^{\rightarrow})$ coordinate-wise, which
certifies that the fetched row is exactly the accepting leaf of that tree on
that sample. Finally, the fetched leaf scores are aggregated back along the
tree axis to obtain the class scores. The remaining forest objects are not used
directly in this fetch-and-containment check, but they are needed to certify
that the committed rows do form a valid hyperbox forest and to connect
inference with the later training proof.

\subsection{Training Proof Pipeline}

The training proof reuses the same committed forest and is organized around one
logical statement per boosting round. For every round $t\in[T]$ and class
$k\in[K]$, the prover exposes a global leaf-assignment vector
\[
\boldsymbol{\ell}^{(t,k)} \in
\{N_{\mathrm{int}},\ldots,N_{\mathrm{tot}}-1\}^{M},
\]
where $\boldsymbol{\ell}^{(t,k)}[r]$ is the global leaf index reached by sample
$r$ in the tree trained for class $k$ at round $t$.

A key structural point is that the family
$(\mathbf{S}^{(0)},\mathbf{S}^{(1)},\ldots,\mathbf{S}^{(T)})$ is not committed
round by round. Instead, all score matrices are arranged along one extra round
axis and committed in a single PCS instance; the same shared commitment pattern
is used for the aligned boosting-state tensors derived from them. The proof
then checks all score-update, softmax, and residual identities against that
shared committed object. Consequently, the proof size does not grow with the
number of boosting rounds; increasing $T$ enlarges the witness tensor and the
prover's work, but not the PCS-level proof artifact.

Logically, the proof of training establishes that for every round $t\in[T]$,
all of the following conditions hold:
\begin{enumerate}
  \item the current boosting state is correct:
  $\mathbf{S}^{(t)}\in\mathbb{F}^{M\times K}$ is the score matrix carried into
  round $t$, and
  $\mathbf{G}^{(t)}\in\mathbb{F}^{M\times K}$ is the pseudo-residual matrix
  derived from $\mathbf{S}^{(t)}$ by the softmax relation;
  \item for every class-specific tree $(t,k)$, the vector
  $\boldsymbol{\ell}^{(t,k)}$ records exactly which leaf each sample reaches in
  the committed tree of that round;
  \item the histogram tensors derived from $\boldsymbol{\ell}^{(t,k)}$ and
  $\mathbf{G}^{(t)}[\cdot,k]$ are correct, first on leaf rows and then after
  propagation to internal nodes through the routing matrices; these tensors lie
  in $\mathbb{F}^{N_{\mathrm{tot}}\times d\times B}$ over the node, feature,
  and bin axes;
  \item from those histograms, the protocol derives the candidate split-score
  tensor
  $\Gamma^{(t,k)}\in\mathbb{F}^{N_{\mathrm{int}}\times d\times B}$ for every
  internal-node / feature / threshold triple, and the split encoded by
  $(\mathbf{E},\boldsymbol{\Theta})$ is certified to optimize that tensor;
  \item the same per-leaf aggregates determine a padded leaf-value matrix
  $\widetilde{\mathbf{W}}^{(t)}\in\mathbb{F}^{N_{\mathrm{tot}}\times K}$ for
  the round-$t$ trees, and these leaf values update the boosting state to the
  next score matrix $\mathbf{S}^{(t+1)}\in\mathbb{F}^{M\times K}$.
\end{enumerate}

Together, these conditions imply that the committed objects encode a valid
plaintext GBDT training trajectory across all rounds. The tree-specific novelty
lies in the middle of this statement: instead of proving a sorting procedure,
the protocol proves the statistics from which split search can be reconstructed
algebraically.

\subsection{Objective Rewriting for Split Search}

Consider one node $u$, one feature $j$, and one threshold bin $b$. Let
$I_u \subseteq [M]$ be the samples reaching node $u$, and let
$L_{u,j,b},R_{u,j,b}$ be the induced left and right partitions. If each sample
$r\in I_u$ carries target value $g_r$, then the residual-sum-of-squares
objective is
\[
\operatorname{RSS}(u,j,b)
=
\sum_{r\in L_{u,j,b}}(g_r-\bar{g}_L)^2
+
\sum_{r\in R_{u,j,b}}(g_r-\bar{g}_R)^2.
\]
Expanding the squares gives
\begin{equation}\label{eq:rss-rewrite}
\operatorname{RSS}(u,j,b)
=
\sum_{r\in I_u} g_r^2
-
\left(
\frac{\left(\sum_{r\in L_{u,j,b}} g_r\right)^2}{|L_{u,j,b}|}
+
\frac{\left(\sum_{r\in R_{u,j,b}} g_r\right)^2}{|R_{u,j,b}|}
\right).
\end{equation}
For fixed node $u$, the first term in \eqref{eq:rss-rewrite} is independent of
the candidate split. Therefore minimizing $\operatorname{RSS}(u,j,b)$ is
equivalent to maximizing
\begin{equation}\label{eq:split-objective}
\Phi(u,j,b)
:=
\frac{\left(\sum_{r\in L_{u,j,b}} g_r\right)^2}{|L_{u,j,b}|}
+
\frac{\left(\sum_{r\in R_{u,j,b}} g_r\right)^2}{|R_{u,j,b}|},
\end{equation}
subject to the validity condition that both children are nonempty.

This rewriting explains why histograms are sufficient proof objects. Once the
prover has committed, for every node $u$, feature $j$, and bin $q$, the sample
count and residual sum associated with that bin, the quantities
$|L_{u,j,b}|$, $|R_{u,j,b}|$,
$\sum_{r\in L_{u,j,b}} g_r$, and $\sum_{r\in R_{u,j,b}} g_r$ become prefix- and
total-sum expressions over the bin axis. The split objective is then reduced to
tensor algebra over committed counts and sums,
which becomes much easier to handle.

\subsection{Domain-Lifted and Interleaved Batching}

The construction repeatedly needs to batch claims that are structurally similar
but do not initially share one proof shape. In this subsection, we describe two
general-purpose batching techniques that are used throughout the construction
to achieve succinctness and modularity. Prior proof systems already rely on
vanishing-polynomial constraints and random-linear-combination batching on a
common domain \cite{aurora19,cq++24}. Our claim is narrower: to the best of our
knowledge, the two batching moves isolated here are new as reusable
proof-composition tools in this setting. The first lifts heterogeneous linear
claims to the same largest domain before batching, and the second interleaves
aligned same-form nonlinear bundles into one pointwise check.

\paragraph{Domain-lifted batch.}
Most linear constraints in the paper are first expressed over some evaluation
domain $\mathbb{H}_n$. After moving all terms to one side, we obtain a virtual polynomial $F(X)$, and the constraint is
reduced to a zero-check over that domain, hence to a divisibility claim of the
form
\[
F(X)=Q(X)\nu_n(X),
\]
where $\nu_n(X)=X^n-1$ is the vanishing polynomial of $\mathbb{H}_n$.

The basic batching problem is therefore the following: suppose we have two such
claims, one over $\mathbb{H}_{n_1}$ and one over $\mathbb{H}_{n_2}$, with
$n_1<n_2$. Since our domain sizes are powers of two, the smaller domain is
nested inside the larger one. We can therefore lift the smaller claim to
$\mathbb{H}_{n_2}$ by multiplying its residual by
$\nu_{n_2}(X)/\nu_{n_1}(X)$:
\[
\widehat{F}_1(X):=F_1(X)\frac{\nu_{n_2}(X)}{\nu_{n_1}(X)}.
\]
Now both claims live over the same ambient domain $\mathbb{H}_{n_2}$, and the
verifier can batch them as
\[
\widehat{F}_1(X)+\eta F_2(X)=Q(X)\nu_{n_2}(X)
\]
for a random challenge $\eta$.

We call this \emph{domain-lifting batching}. It lets us merge linear claims
that are initially written over different domains by first rewriting the
smaller one over the larger domain and then applying an ordinary same-domain
batch.

\paragraph{Interleaving batching.}
The second batching problem concerns same-form nonlinear checks that are
\emph{local}: validity at one position depends only on the witness values at
that position, not on any other entry of the vector. Typical examples in our
setting include lookup-style checks against one common table and range checks
against one public range table. If each local instance is proved per PCS, the
proof size
scales with the number of instances even though the check itself is always the
same.

The basic case is two vectors of the same length. Suppose
$\mathbf{f},\mathbf{g}\in\mathbb{F}^n$ must both satisfy the same local
constraint, say $\varphi(x)=0$ entrywise. Then we can interleave them into one
larger virtual vector
\[
(\mathbf{f}[0],\mathbf{g}[0],\mathbf{f}[1],\mathbf{g}[1],\ldots),
\]
and check the same local constraint once on that merged vector. In polynomial
form, if $f(X)$ and $g(X)$ are the encodings of $\mathbf{f}$ and $\mathbf{g}$
over $\mathbb{H}_n$, the merged object is
\[
\operatorname{Int}_n(f,g)(X)
:=
\operatorname{even}_n(X)\,f(X)
+ \operatorname{odd}_n(X)\,g(\omega_{2n}^{-1}X),
\]
where
\[
\operatorname{even}_n(X):=\frac{X^n+1}{2},
\qquad
\operatorname{odd}_n(X):=\frac{1-X^n}{2}.
\]
These public sparse selectors place $\mathbf{f}$ on the even positions and
$\mathbf{g}$ on the odd positions, so every position of the larger vector still
corresponds to one concrete original entry.

If the two vectors have different lengths, we first lift the smaller one to the
larger domain by the substitution $X\mapsto X^r$, exactly as in \cref{lem:row-replication}.
The important point is that the merged object is still virtual: it is assembled
from existing commitments, public sparse masks, and domain-lift substitutions,
so no new base commitment is required.
